% 5.3 空间群表示的标记
\section{空间群表示的标记}

在表 5.8 中我们给出了用于 ${}^{H}\mathbf{G}^{\mathbf{k}}$ 或 $\overline{\mathbf{G}}^{\mathbf{k}*}$ 的表示的标记。我们使用表 2.2 中点群表示所用的两种记号的推广，即 Mulliken (1933) 的记号（$A$、$B$、$E$、$T$ 等）以及 Koster, Dimmock, Wheeler, and Statz (1963) 的 $\Gamma_{1}, \Gamma_{2}, \Gamma_{3}$ 等记号。对许多空间群，${}^{H}\mathbf{G}^{\mathbf{k}}$ 或 $\overline{\mathbf{G}}^{\mathbf{k}*}$ 同构于 32 个点群之一，此时我们直接使用表 2.2 中该点群的标记。当 ${}^{H}\mathbf{G}^{\mathbf{k}}$ 或 $\overline{\mathbf{G}}^{\mathbf{k}*}$ 不同构于任何点群时，记号作如下推广（这里“实”指所有特征标都是实的，“复”指部分特征标是复的）：

\begin{center}
\begin{tabular}{ccl}
\toprule
表示 & 维数 & 标记\\
\midrule
实 & 1 & $A$ 或 $B$\\
复 & 1 & ${}^{1}E,\ {}^{2}E$（共轭对）\\
实 & 2 & $E$\\
复 & 2 & ${}^{1}F,\ {}^{2}F$（共轭对）\\
实 & 3 & $T$\\
复 & 3 & ${}^{1}H,\ {}^{2}H$（共轭对）\\
实 & 4 & $F$\\
复 & 4 & ${}^{1}J,\ {}^{2}J$（共轭对）\\
实 & 6 & $H$\\
\bottomrule
\end{tabular}
\end{center}

下标 1、2、3 等，下标 g 与 u，以及上标撇与双撇，用于区分同维数的表示，方式与点群标记中类似。若 ${}^{H}\mathbf{G}^{\mathbf{k}}$ 是 $(\{E \mid \mathbf{0}\} \oplus \{I \mid \mathbf{0}\})$ 与某个群 $\mathbf{G}$ 的直积，则用 $R_{g}$ 与 $R_{u}$ 表示由 $\mathbf{G}$ 的表示 $R$ 导出的两个表示：在 $R_{g}$ 中 $\chi(\{I \mid \mathbf{0}\})$ 为 $+\chi(\{E \mid \mathbf{0}\})$，在 $R_{u}$ 中 $\chi(\{I \mid \mathbf{0}\})$ 为 $-\chi(\{E \mid \mathbf{0}\})$。撇与双撇用于区分仅仅在一个重要反射面的特征标符号上不同的两个表示 $R'_{i}$ 与 $R''_{i}$。在 $\Gamma$ 记法中，表示标记为 $\Gamma_{1}, \Gamma_{2}, \Gamma_{3}$ 等；上标 $+$ 与 $-$ 通常用于直积群的表示，用法与上面的 g 和 u 类似（Miller and Love (1967) 除外）。在对称点上（$\Gamma$ 以外的点）或对称线上，$\Gamma$ 换成表示该点或线的字母。例如，在 $P4_{3}32\ (O^{6})$ 的 $M$ 点，与抽象群 $\mathbf{G}^{10}_{16}$ 的表示 $R_{5}, R_{6}, R_{7}, R_{8}, R_{9}$ 认同的空间群表示，按表 5.8 中 $\mathbf{G}^{10}_{16}$ 下的条目“a”，或者标记为 ${}^{2}E_{1}, {}^{1}E_{1}, {}^{1}E_{2}, {}^{2}E_{2}$ 与 $E$，或者标记为 $M_{1}, M_{2}, M_{3}, M_{4}$ 与 $M_{5}$，其中 $\Gamma_{i}$ 已换成 $M_{i}$。

在这两种记法中都不可能完全合乎逻辑地定义一套标记。也难于综合研究过个别空间群或少量空间群的各位作者所用的所有标记。

\begin{figure}[H]
\centering
\includegraphics[width=0.86\textwidth]{c5_t58_p01.jpg}
\caption*{\textbf{表 5.8}\ 空间群表示的标记（原书第 390 页起扫描；第 1 页，共 7 页）。}
\end{figure}
\begin{figure}[H]
\centering
\includegraphics[width=0.86\textwidth]{c5_t58_p02.jpg}
\caption*{\textbf{表 5.8}（第 2 页）。}
\end{figure}
\begin{figure}[H]
\centering
\includegraphics[width=0.86\textwidth]{c5_t58_p03.jpg}
\caption*{\textbf{表 5.8}（第 3 页）。}
\end{figure}
\begin{figure}[H]
\centering
\includegraphics[width=0.86\textwidth]{c5_t58_p04.jpg}
\caption*{\textbf{表 5.8}（第 4 页）。}
\end{figure}
\begin{figure}[H]
\centering
\includegraphics[width=0.86\textwidth]{c5_t58_p05.jpg}
\caption*{\textbf{表 5.8}（第 5 页）。}
\end{figure}
\begin{figure}[H]
\centering
\includegraphics[width=0.86\textwidth]{c5_t58_p06.jpg}
\caption*{\textbf{表 5.8}（第 6 页）。}
\end{figure}
\begin{figure}[H]
\centering
\includegraphics[width=0.86\textwidth]{c5_t58_p07.jpg}
\caption*{\textbf{表 5.8}（第 7 页）。}
\end{figure}
